% Generated draft: requires mathematical and source review.
\documentclass[11pt]{article}
\usepackage[margin=25mm]{geometry}
\usepackage{fontspec}
\setmainfont{Latin Modern Roman}
\usepackage{amsmath,amssymb,mathtools,unicode-math}
\setmathfont{Latin Modern Math}
\usepackage{enumitem,needspace,xurl,hyperref}
\setlength{\parindent}{0pt}
\setlength{\parskip}{5pt}
\setlength{\emergencystretch}{3em}
\newcommand{\N}{\mathbb{N}}
\newcommand{\Z}{\mathbb{Z}}
\newcommand{\Q}{\mathbb{Q}}
\newcommand{\R}{\mathbb{R}}
\newcommand{\C}{\mathbb{C}}
\newcommand{\F}{\mathbb{F}}
\newcommand{\A}{\mathbb{A}}
\newcommand{\PP}{\mathbb{P}}
\newcommand{\E}{\mathbb{E}}
\newcommand{\eps}{\varepsilon}
\newcommand{\dd}{\,\mathrm{d}}
\begin{document}
\section*{UnsolvedMath ID: 30000376}
\textbf{Source problem code:} catalogue entry
\subsection{Definitions and mathematical statement}
\textbf{Definability-complexity question.} What is the least L\'evy-hierarchy complexity of a parameter-free formula \(\Phi(x,y)\) for which there is a set-forcing extension in which \(\Phi\) defines a well-order of the reals over \(\langle H(\omega_2),\in\rangle\), or over a specified natural expansion such as \(\langle H(\omega_2),\in,\mathrm{NS}_{\omega_1}\rangle\)? The base structure and expanded structure are separate optimization problems; adding a predicate changes definability complexity.\par In a forcing extension \(V[G]\), let \(H(\omega_2)=\{a:|\operatorname{tc}(\{a\})|<\omega_2\}\), with hereditary size and the cardinals computed in that extension. Reals may be coded as subsets of \(\omega\). The unary predicate \(\mathrm{NS}_{\omega_1}\) names the nonstationary subsets of \(\omega_1\): \(A\subseteq\omega_1\) is nonstationary if it is disjoint from a closed unbounded subset of \(\omega_1\). This is a predicate on the structure, not an assumption that the ideal itself belongs to \(H(\omega_2)\). Parameter-free means no set, real, or ordinal parameters are allowed in \(\Phi\), beyond its displayed free variables; the fixed language predicates remain available.\par For a fixed permitted ground model and forcing, the relation is\[x\prec y\quad\Longleftrightarrow\quad\langle H(\omega_2)^{V[G]},\in\rangle\models\Phi(x,y),\] with the additional predicate interpreted in the expansion when appropriate. It must be a well-order of all reals in \(V[G]\). The L\'evy hierarchy starts with bounded formulas \(\Delta_0\) and counts unbounded quantifier alternations to obtain \(\Sigma_n\) and \(\Pi_n\), relative to the chosen language; a \(\Delta_n\) definability claim requires equivalent definitions of both kinds. Thus optimizing \(\Sigma_n\), \(\Pi_n\), or \(\Delta_n\) must be distinguished, and no claim about projective complexity is silently substituted.\par The source does not fix the starting theory or ground-model hypotheses, permitted forcings and preservation conditions, or whether one formula must work uniformly over all allowed ground models. Existence over some ground model, existence over each model in a class, and existence over a specified model are different assertions. It also does not enumerate the allowed natural expansions; arbitrary added predicates coding a chosen well-order would change the task. These missing quantifiers prevent a unique numerical optimum from being extracted.\par The supplied triage checked 2026-08-21 labels the item partially solved. It reports forcing parameter-free definable well-orders of \(H(\omega_2)\) under explicit cardinal-arithmetic hypotheses, while describing the general optimal complexity for well-ordering the reals in the listed structures as unsettled. Those hypotheses and exact complexity levels are not reproduced, so the result for all of \(H(\omega_2)\) is not silently used as an optimality theorem for this question. The dated assessment requires independent verification.
\subsection{Short English statement}
Find the lowest definability level at which forcing can produce a parameter-free well-order of all reals, separately in the membership structure and its specified expansions. The ground-model and forcing quantifiers, preservation requirements, and choice among \(\Sigma\), \(\Pi\), and \(\Delta\) complexity are essential missing inputs.
\subsection{Sources}
\begin{itemize}\item Local UnsolvedMath registry, permanent ID 30000376, code catalogue entry; \emph{Definability Complexity of Well-Orderings of the Reals}, Oberwolfach Reports Open Problems.\item \emph{Set Theory (2006)}, Oberwolfach Reports, DOI 10.4171/owr/2005/55, pp. 3121--3174. \url{https://doi.org/10.4171/owr/2005/55}. Proposed year and year tag: 2005; no proposer named.\item Supplied literature label \emph{original OWR report}: \url{https://doi.org/10.4171/owr/2005/55}.\item Supplied literature label \emph{definable well-orders of \(H(\omega_2)\)}: \url{https://doi.org/10.2178/jsl/7704030}.\item Literature triage checked 2026-08-21, status \texttt{partially\_solved}; record updated 2026-09-16. The \texttt{corrected\_verified} field reports checking both the base structure and nonstationary-ideal expansion in the official report, not independently checked here.\end{itemize}
\end{document}
